
\section{Proof of Theorem~\ref{thm:main} and Experimental Validation}\label{app:proof}

\noindent {\bf Proof sketch.}

\noindent\textbf{Step 1: Low-degree nodes are weakly coupled.}
A node with degree $k_i \leq d$ participates in at most $d$ edges. Moving it between communities changes modularity by at most $O(d/m)$: the adjacency term shifts by at most $d/m$, and the null-model penalty shifts by $O(k_i \cdot K_r/(2m)^2)$ where $K_r$ is the total degree of the destination community. In the sparse regime where $m = \Theta(n)$, this is $O(1/n)$, so the node's community assignment barely affects the objective.

\medskip
\noindent\textbf{Step 2: Weakly coupled nodes can be independently reassigned.}
Reassigning one low-degree node perturbs another's sensitivity by only $O(1/n^2)$. By selecting a large independent set among the low-degree nodes, which is straightforward since these nodes have bounded degree, we obtain $\Theta(n)$ nodes whose reassignments are approximately independent. Each can be placed in at least two communities without meaningful modularity loss, yielding $2^{\Theta(n)}$ distinct near-optimal partitions.

\medskip
\noindent\textbf{Step 3: Sparsity makes this severe.}
When $\kbar < 3$, standard results on Poisson and power-law degree distributions guarantee that $n_{\mathrm{\le d}} = \Theta(n)$. This is the regime of knowledge graphs in GraphRAG. In contrast, for dense graphs with large $\kbar$, $n_{\mathrm{\le d}}/n \to 0$ and most nodes are strongly coupled to their communities, collapsing the degeneracy.\\

\noindent {\bf Full proof.}

\begin{proof}
We first remind modularity in its standard community-sum form:
\[
Q(\sigma) = \sum_{c} \left[\frac{e_c}{m} - \left(\frac{K_c}{2m}\right)^{\!2}\right],
\]
where $e_c$ counts edges internal to community $c$ and $K_c = \sum_{j \in c} k_j$ is its total degree. Moving node $i$ from community $s$ to community $r$ produces a new partition $\sigma^{i \to r}$. We write $d_i^{(c)}$ for the number of neighbors of $i$ in community $c$.
We also define the \emph{modularity sensitivity} of node $i$ under partition $\sigma$ is $\Delta_i(\sigma) = \max_{r \neq \sigma_i} |Q(\sigma) - Q(\sigma^{i \to r})|$, where $\sigma^{i \to r}$ moves $i$ to community $r$ with all other assignments fixed. Node $i$ is \emph{$\varepsilon$-weakly coupled} if $\Delta_i(\sigma) < \varepsilon$.

Next, we give two lemmas that will help with the proof

\begin{lemma}[Moving a low-degree node]\label{lem:one}
For any partition $\sigma$ and any node $i$ with degree $k_i \leq d$, moving $i$ to any other community changes modularity by at most
\[
\Delta_i(\sigma) \;\leq\; \frac{2k_i}{m} + \frac{k_i^2}{2m^2} \;=\; O\!\left(\frac{1}{n}\right).
\]
\end{lemma}

\begin{proof}
Moving $i$ from community $s$ to $r$ affects only those two terms in the sum. There are two contributions:

\medskip
\noindent\textit{Edge term.} Community $s$ loses $d_i^{(s)}$ internal edges; $r$ gains $d_i^{(r)}$. Both are at most $k_i$, so
\[
|\Delta_{\mathrm{edge}}| = \frac{|d_i^{(r)} - d_i^{(s)}|}{m} \leq \frac{k_i}{m}.
\]

\noindent\textit{Degree-penalty term.} The total degree of $s$ drops by $k_i$ and that of $r$ rises by $k_i$. Expanding the squares and simplifying:
\[
|\Delta_{\mathrm{penalty}}| = \frac{|k_i[2(K_r - K_s) + 2k_i]|}{(2m)^2} \leq \frac{k_i}{m} + \frac{k_i^2}{2m^2},
\]
using $|K_r - K_s| \leq 2m$.

\medskip
\noindent Adding the two and substituting $k_i \leq d$ and $m = n\kbar/2 = \Theta(n)$ gives $\Delta_i \leq 2d/m + d^2/(2m^2)$. Both terms are $O(1/n)$.
\end{proof}

\begin{lemma}[Non-adjacent low-degree nodes]\label{lem:two}
Let $i,j$ be non-adjacent nodes with degrees $\leq d$. Reassigning $j$ changes node $i$'s sensitivity by at most
\[
|\Delta_i(\sigma) - \Delta_i(\sigma')| \;\leq\; \frac{2d^2}{(2m)^2} \;=\; O\!\left(\frac{1}{n^2}\right).
\]
\end{lemma}

\begin{proof}
Since $i$ and $j$ share no edge, reassigning $j$ does not change $i$'s neighbor counts $d_i^{(c)}$ in any community, so the edge term for $i$ is unaffected. The only effect is that $j$'s move shifts the total degrees $K_r, K_s$ by $\pm k_j \leq d$, which perturbs $i$'s penalty term by at most $k_i \cdot 2d/(2m)^2 \leq 2d^2/(2m)^2 = O(1/n^2)$.
\end{proof}

\noindent Now, we give the proof of Theorem~\ref{thm:main}. The argument has three parts.

\medskip
\noindent\textbf{1. Find a large independent set of weakly coupled nodes.}
By Lemma~\ref{lem:one}, every node with $k_i \leq d$ has sensitivity $O(1/n)$, so for $\varepsilon > d(2+\kbar)/(2m)$ it is $\varepsilon$-weakly coupled. Let $n_{\le d}$ denote the number of such nodes. The subgraph they induce has maximum degree $d$, so a greedy algorithm yields an independent set $S$ of size
\[
|S| \;\geq\; \frac{n_{\le d}}{d+1}.
\]

\medskip
\noindent\textbf{2. Every combination of reassignments stays near-optimal.}
Start from an optimal partition $\sigma^*$. Each node $i \in S$ can stay or move to an alternative community: two choices per node. For any subset $T \subseteq S$ of nodes that are reassigned, the total modularity loss is bounded by summing the individual sensitivities plus pairwise cross-terms (using Lemma~\ref{lem:two}):
\[
|Q^* - Q(\sigma_T)| \;\leq\; \underbrace{\sum_{i \in T} \Delta_i(\sigma^*)}_{\text{individual}} \;+\; \underbrace{\sum_{\{i,j\} \subseteq T} \frac{2d^2}{(2m)^2}}_{\text{cross-terms}}.
\]
Bounding each sum with $|T| \leq |S| \leq n_{\le d}/(d+1)$ and $m = n\kbar/2$:
\begin{align*}
\text{Individual sum} &\leq \frac{n_{\le d}}{d+1} \cdot \frac{d(2+\kbar)}{n\kbar} \leq \frac{d(2+\kbar)}{(d+1)\kbar} \eqqcolon C_1, \\[6pt]
\text{Cross-term sum} &\leq \frac{1}{2}\left(\frac{n_{\le d}}{d+1}\right)^{2} \cdot \frac{2d^2}{n^2\kbar^2} \leq \frac{d^2}{(d+1)^2\kbar^2} \eqqcolon C_2.
\end{align*}
Both $C_1$ and $C_2$ are constants (independent of $n$), so setting $\varepsilon' = C_1 + C_2$ keeps every such partition within $\varepsilon'$ of $Q^*$.

\medskip
\noindent\textbf{3. Count and conclude.}
Each of the $|S|$ nodes has two valid assignments, giving $2^{|S|}$ structurally distinct near-optimal partitions:
\[
\mathcal{D}(\varepsilon') \;\geq\; 2^{|S|} \;\geq\; 2^{\,n_{\le d}/(d+1)}.
\]
In the sparse regime ($\kbar < 3$), standard degree distributions guarantee $n_{\le d} = \Theta(n)$. For example, under a Poisson distribution with $\kbar = 2$ and $d = 2$, the fraction of low-degree nodes is $5e^{-2} \approx 0.68$; under a power law with exponent $\gamma = 2.5$, degree-1 nodes alone comprise ${\sim}74\%$ of the graph. Therefore $\mathcal{D}(\varepsilon') \geq 2^{\Theta(n)}$.
\end{proof}

\subsection{Empirical Analysis of Leiden Instability}
To complement Theorem~\ref{thm:main}, we conducted a multi-seed stability analysis (10 runs with different random seeds) of Leiden across 3 datasets (GPT-3.5 setup). 
Across runs, the cluster statistics show noticeable variability. For example, on the Podcast dataset, the number of clusters ranges from 2550 to 2584, while the largest cluster size ranges between 780 and 1069 nodes. The average node-level Adjusted Rand Index (ARI) is 0.94, indicating around 6\% disagreement in node-pair assignments.
While moderate, this level of variation is significant in GraphRAG as community assignments directly define retrieval units and contexts for the LLM: the Jaccard similarity of retrieved nodes across seeds is ~0.73, on average.
These findings provide direct empirical evidence of retrieval-level instability in Leiden clustering and further support Theorem~\ref{thm:main}.


\section{Splitting Oversized Components}

\subsection{SPLIT: Splitting Large Connected Components} \label{sec:split_cc}

When a connected component $R$ exceeds the maximum cluster size $M$, it is split into smaller clusters using the \textsc{Split} procedure. The algorithm (Algorithm~\ref{alg:split_cc}) grows each cluster greedily from a seed node with the highest degree (line~\ref{line:pick_seed}), adding neighboring nodes that maximize internal connectivity (line~\ref{line:add_best_node}) until the cluster reaches size $M$. Each completed cluster is added to the cluster set (line~\ref{line:update_cluster}) and removed from the remaining node set (line~\ref{line:remove_set}). The output is a set of size-constrained clusters that preserve internal connectivity within $C$.

\begin{algorithm}[!h]
\caption{\textsc{Split} (G, R, M)}
\label{alg:split_cc}
\begin{algorithmic}[1]
\REQUIRE Graph $G=(V,E)$, connected comp. $R$, max size $M$
\ENSURE List of split clusters $\mathcal{C}$

%\STATE $N \gets$ set of nodes in the component
\STATE Initialize cluster set $\mathcal{C} \gets \emptyset$
\WHILE{$R \neq \emptyset$}
    \STATE Pick seed node $s \in R$ with highest degree \label{line:pick_seed}
    \STATE $S \gets \{s\}$, $frontier \gets N(s) \cap R$
    
    \WHILE{$|S| < M$ \textbf{and} $frontier \neq \emptyset$}
        \STATE Pick $v \in frontier$ maximizing $N(v) \cap S$ \label{line:add_best_node}
        \STATE $S \gets S \cup v$
        \STATE $frontier \gets N(v) \cap (R \setminus S)$
    \ENDWHILE
        \STATE $\mathcal{C} \gets \mathcal{C} \cup S $  \label{line:update_cluster}
%    \STATE Add $S$ to $clusters$
    \STATE Remove $S$ from $R$ \label{line:remove_set}
\ENDWHILE

\RETURN $\mathcal{C}$
\end{algorithmic}
\end{algorithm}


\FloatBarrier


\subsection{SPLIT-2HOP: Splitting Large 2-Hop Connected Components}
\label{sec:split2hop}

When a 2-hop connected component $H$ exceeds the maximum cluster size $M$, it is partitioned into smaller clusters using the \textsc{Split-2HOP} procedure (Algorithm~\ref{alg:split2hop}). The algorithm first computes the anchor set as the union of neighbors of all 2-hop nodes and precomputes anchor neighborhoods for each node (lines~\ref{line:anchors}--\ref{line:Au}).  
Clusters are grown greedily by selecting a seed node with the largest anchor set (line~\ref{line:seed}) and initializing a cluster with this seed (line~\ref{line:initS}). Candidate nodes that share at least one anchor with the seed are added to the frontier (line~\ref{line:frontier}). The cluster is expanded by repeatedly selecting the frontier node with the highest overlap in anchor connectivity with the current cluster (lines~\ref{line:innerwhile}--\ref{line:addv}), until the cluster reaches size $M$ or no eligible nodes remain.  
After cluster growth, anchor nodes connected to at least two nodes in the cluster are selected and included (line~\ref{line:selectanchors}), and the resulting cluster is added to the output set (line~\ref{line:addcluster}). The assigned nodes are then removed from the remaining 2-hop groups set (line~\ref{line:removeS}). This process repeats until all nodes in the 2-hop groups are assigned.


\begin{algorithm}[!htb]
\caption{\textsc{Split-2hop} (G, H, M)}
\label{alg:split2hop}
\begin{algorithmic}[1]
\REQUIRE Graph $G=(V,E)$, 2-hop groups $H$, max size $M$
\ENSURE List of split clusters $\mathcal{C}$

\STATE Initialize cluster set $\mathcal{C} \gets \emptyset$ \label{line:initC}

\STATE Compute anchors $A \gets \bigcup_{u \in H} N(u)$ \label{line:anchors}
\STATE Precompute $A_u \gets N(u) \cap A$ for $u \in H$ \label{line:Au}

\WHILE{$H \neq \emptyset$} \label{line:outerwhile}
    \STATE Pick $seed \in H$ with max $|A_{seed}|$ \label{line:seed}
    \STATE $S \gets \{seed\}$ \label{line:initS}
    \STATE $frontier \gets \{v \in H\setminus\{seed\} \mid A_v \cap A_{seed} \neq \emptyset\}$ \label{line:frontier}

    \WHILE{$|S| < M$ and $frontier \neq \emptyset$} \label{line:innerwhile}
        \STATE Pick $v \in frontier$ maximizing $\sum_{u \in S} |A_v \cap A_u|$ \label{line:pickv}
        \STATE $S \gets S \cup \{v\}$\;
        $frontier \gets (frontier \cup \{u \in H \mid A_u \cap A_v \neq \emptyset\}) \setminus S$ \label{line:addv}
    \ENDWHILE
    
\tcp{Select anchors linked to $\ge 2$ nodes in $S$} 
    \STATE $A' \gets \{v \in A s.t. N(v) \cap S \ge 2\}$ \label{line:selectanchors}
    \STATE $\mathcal{C} \gets \mathcal{C} \cup \{S \cup A'\}$ \label{line:addcluster}
    \STATE $H \gets H \setminus S$ \label{line:removeS}
\ENDWHILE

\RETURN $\mathcal{C}$
\end{algorithmic}
\end{algorithm}



\section{Alternative Community Partition Methods} \label{sec:altcomdet}
To our knowledge, prior GraphRAG studies do not evaluate alternative community detection methods, making Leiden the de facto comparison baseline. We also explored hierarchical block structures \cite{peixoto14} and non-modularity-based approaches such as $k$-truss, but found them consistently weaker than Leiden and therefore omitted them from the experimental results (additional discussion on $k$-truss is provided in Section~\ref{sec:proposed_heuristics}). For example, on the Semiconductor dataset (GPT-3.5 setup), Leiden outperforms $k$-truss in $\sim$82--85\% of comparisons, and $k$-core outperforms $k$-truss in $\sim$90--93\% of comparisons.


\section{Evaluation Criteria Details} 
\subsection{Question Generation} \label{sec:question_generation}
Following methodology described in \cite{edge24}, we generated 125 sensemaking questions using \five.
Consistent with their approach, we first prompted \five to create personas of hypothetical users for each corpus and then generated 5 tasks per user.  
For each user-task pair, we use \five to generate high-level questions that require understanding of the entire corpus without relying on low-level fact retrieval.  
This process ensures that the resulting questions assess comprehensive, corpus-wide reasoning, aligned with prior work.

\subsection{Evaluation Approach} \label{sec:evaluation_details}
We evaluate models using a head-to-head framework, where LLMs compare pairs of answers to identify a winner, loser, or tie. To enhance reliability, multiple evaluators are used, and repeated assessments are performed. The final outcome is determined via majority voting, making this approach suitable for global sensemaking tasks without gold-standard references.
 Each comparison proceeds as follows:
\begin{enumerate}
\item Generate answers to the same query using two approaches.
\item Randomize answer order and present the pair to the evaluator.
\item Each of five independent LLM evaluators (\five, Gemini 3 Pro Preview, Gemini 2.5 Pro, Qwen3 Next 80B, and DeepSeek v3.2) assesses the same answer pair three times.
\item Apply majority voting within each evaluator to determine their decision.
\item Apply a second majority vote across evaluators to determine the final judgment.
\end{enumerate}
This procedure ensures robust, repeatable head-to-head comparisons across methods.


%cut starting here in the camera-ready

\section{Statistical Analysis (p-values)} \label{sec:p_values_appendix}
To assess the significance of observed differences, we follow a procedure similar to Edge et al.~\cite{edge24}. 
For each dataset, evaluator, and metric, we assign scores for each head-to-head comparison: the winning method receives a score of 100, the losing method receives 0, and in the event of a tie, both methods receive 50.  These scores are then averaged over all evaluators, repeated runs, and questions, as detailed in Section~\ref{sec:evaluation_method}.
Following~\cite{edge24}, we use non-parametric Wilcoxon signed-rank tests to evaluate pairwise performance differences between methods. 
Holm-Bonferroni adjustment is applied to correct for multiple comparisons.


\begin{table}[!t]
\centering
\caption{Wilcoxon signed-rank test p-values for comparisons across datasets and baselines on Diversity (\threepointfive). Significant values are highlighted in \textbf{bold} for $p<0.005$.}
\label{tab:p_values_diversity}
\resizebox{.47\textwidth}{!}{%
\begin{tabular}{l c c | c c | c c}
\hline
\textbf{Condition} & \multicolumn{2}{c|}{\podcast} & \multicolumn{2}{c|}{\news} & \multicolumn{2}{c}{\semi} \\
 & C2 & C3 & C2 & C3 & C2 & C3 \\
\hline
\KHC L1   
& 0.073 & \textbf{<0.001} 
& 0.008 & \textbf{<0.001} 
& 0.083 & 0.083 \\

\KHC LF   
& \textbf{<0.001} & \textbf{<0.001} 
& 0.071 & \textbf{<0.001} 
& 0.008 & 0.007 \\

\MHC L1   
& \textbf{<0.001} & \textbf{<0.001} 
& \textbf{<0.001} & 0.007 
& \textbf{<0.001} & 0.982 \\

\MHC LF   
& \textbf{<0.001} & \textbf{<0.001} 
& \textbf{<0.001} & \textbf{<0.001} 
& 0.083 & \textbf{<0.001} \\

\MRC L1   
& 0.982 & \textbf{<0.001} 
& \textbf{<0.001} & 0.071 
& \textbf{<0.001} & \textbf{<0.001} \\

\MRC LF   
& 0.083 & 0.393 
& 0.083 & \textbf{<0.001} 
& \textbf{<0.001} & \textbf{<0.001} \\

\hline
\end{tabular}%
}
\end{table}




%\clearpage


\subsection{Head-to-Head Comparison on Leiden C2 and Different \MHC Levels}

Table~\ref{tab:gpt35_msft_mhc_h2h} shows the head-to-head win rates (\%) for \threepointfive on \ms data, comparing the Leiden C2 baseline against various \MHC level variants. The two subtables report metrics for (a) Comprehensiveness and (b) Diversity.  

Here, the \MHC variants correspond to different k-core levels within the graph hierarchy: LF represents the leaf-level communities, LF-2 corresponds to two levels above the leaf, LF-4 to four levels above the leaf, and LF-6 to six levels above the leaf. We observe that higher k-core levels consistently achieve better win rates against C2, particularly in Comprehensiveness, indicating that the more central nodes capture global context more effectively. 
This justifies focusing on leaf-level (LF) and the level immediately above leaf (L1) for all of our evaluations, as they provide better performance.

\begin{table}[h!]
\centering
\caption{\threepointfive Head-to-Head Win Rates (\%) using Leiden C2 and \MHC Level Variants on \ms data}
\label{tab:gpt35_msft_mhc_h2h}
\setlength{\tabcolsep}{2pt} % reduce column padding

\begin{subtable}[t]{\columnwidth}
\centering
\caption{Comprehensiveness}
\begin{tabular}{|l|c|c|c|c|c|}
\hline
\textbf{} & Leiden C2 & \MHC LF-6 & \MHC LF-4 & \MHC LF-2 & \MHC LF \\
\hline
Leiden C2 & \heatmap{50} & \heatmap{56} & \heatmap{51} & \heatmap{48} & \heatmap{39} \\
\MHC LF-6 & \heatmap{44} & \heatmap{50} & \heatmap{45} & \heatmap{43} & \heatmap{33} \\
\MHC LF-4 & \heatmap{49} & \heatmap{55} & \heatmap{50} & \heatmap{46} & \heatmap{40} \\
\MHC LF-2 & \heatmap{52} & \heatmap{57} & \heatmap{54} & \heatmap{50} & \heatmap{47} \\
\MHC LF   & \heatmap{61} & \heatmap{67} & \heatmap{60} & \heatmap{53} & \heatmap{50} \\
\hline
\end{tabular}
\end{subtable}

\vspace{0.5em} % small space between subtables

\begin{subtable}[t]{\columnwidth}
\centering
\caption{Diversity}
\begin{tabular}{|l|c|c|c|c|c|}
\hline
\textbf{} & Leiden C2 & \MHC LF-6 & \MHC LF-4 & \MHC LF-2 & \MHC LF \\
\hline
Leiden C2 & \heatmap{50} & \heatmap{62} & \heatmap{54} & \heatmap{46} & \heatmap{39} \\
\MHC LF-6 & \heatmap{38} & \heatmap{50} & \heatmap{42} & \heatmap{38} & \heatmap{36} \\
\MHC LF-4 & \heatmap{46} & \heatmap{58} & \heatmap{50} & \heatmap{49} & \heatmap{41} \\
\MHC LF-2 & \heatmap{54} & \heatmap{62} & \heatmap{51} & \heatmap{50} & \heatmap{51} \\
\MHC LF   & \heatmap{61} & \heatmap{64} & \heatmap{59} & \heatmap{49} & \heatmap{50} \\
\hline
\end{tabular}
\end{subtable}

\end{table}






\section{Additional Results}

Here, we present additional results from \threepointfive, \four, and \five, including the results from Leiden C0, C1, and other metrics such as Empowerment and Directness.



\subsection{\four Results on Post-cutoff Data}
Table~\ref{tab:gpt4} reports head-to-head win rates for comprehensiveness and diversity under the same experimental conditions as in \threepointfive.
Across datasets and configurations, our heuristics still achieve higher win rates than Leiden C2/C3 in a majority of comparisons, averaging approximately \textbf{45--55\%}, particularly on \semi.
Among our methods, {\bf \MHC LF still achieves the strongest overall performance}, with average win rates of approximately \textbf{50--52\%} across datasets and metrics, followed by \KHC LF and \MRC LF at \textbf{48--50\%}. 
Dataset-wise, gains are most pronounced on \semi, where \MHC LF reaches up to \textbf{64\%} wins for diversity (C2) and maintains \textbf{58\%} against C3, while comprehensiveness peaks at \textbf{52\%}. 
For \podcast and \news, win rates remain closer to parity, typically fluctuating between \textbf{45--50\%} across heuristics and Leiden levels.

Comparing Leiden resolutions, differences between C2 and C3 are smaller than for \threepointfive, with average win-rate gaps generally within \textbf{2--4\%}. 
Averaged across datasets and metrics, leaf-level (LF) variants continue to outperform L1, but by a narrower margin of \textbf{2--5} percentage points.
Overall, \four narrows the gap between $k$-core heuristics and Leiden GraphRAG, with \textbf{\MHC LF} remaining the most robust, especially on \semi, and generally matching or slightly surpassing Leiden baselines despite more ties.

Note that performance on the \podcast is less reliable due to the limited number of post-cutoff documents. With only 13 documents available, the resulting graphs are sparse, and our heuristics often show minimal improvement, performing similarly to C2/C3 configurations.




\begin{table}[!ht]
\centering
\caption{\five full: Head-to-head win rates (\%), for comprehensiveness and diversity metrics. C3 is the Leiden community level from Edge et al.~\cite{edge24}. LF indicates leaf-level communities, and L1 indicates the level immediately above the leaf.}
\label{tab:gpt5_full_c3}
\resizebox{.47\textwidth}{!}{%
\begin{tabular}{|r|ccc|ccc|ccc|}
\hline
\five
& \multicolumn{3}{c|}{\podcast} 
& \multicolumn{3}{c|}{\news} 
& \multicolumn{3}{c|}{\ms} \\
\cline{2-10}
results for
& \multicolumn{3}{c|}{C3}
& \multicolumn{3}{c|}{C3}
& \multicolumn{3}{c|}{C3} \\
\cline{2-10}
{\bf Comprehensiveness}
& Win & Loss & Tie
& Win & Loss & Tie
& Win & Loss & Tie \\
\hline
\KHC L1   
& \heatmap{41} & \heatmap{41} & \heatmap{18}
& \heatmap{52} & \heatmap{39} & \heatmap{9}
& \heatmap{38} & \heatmap{43} & \heatmap{19} \\

\KHC LF   
& \heatmap{43} & \heatmap{43} & \heatmap{14}
& \heatmap{47} & \heatmap{45} & \heatmap{8}
& \heatmap{42} & \heatmap{48} & \heatmap{10} \\

\MHC L1  
& \heatmap{48} & \heatmap{48} & \heatmap{4}
& \heatmap{49} & \heatmap{37} & \heatmap{14}
& \heatmap{48} & \heatmap{42} & \heatmap{10} \\

\MHC LF  
& \heatmap{50} & \heatmap{47} & \heatmap{3}
& \heatmap{47} & \heatmap{35} & \heatmap{18}
& \heatmap{47} & \heatmap{48} & \heatmap{5} \\

\MRC L1 
& \heatmap{45} & \heatmap{45} & \heatmap{10}
& \heatmap{32} & \heatmap{48} & \heatmap{20}
& \heatmap{45} & \heatmap{40} & \heatmap{15} \\

\MRC LF 
& \heatmap{50} & \heatmap{37} & \heatmap{13}
& \heatmap{36} & \heatmap{46} & \heatmap{18}
& \heatmap{45} & \heatmap{40} & \heatmap{15} \\
\hline

{\bf Diversity}
& Win & Loss & Tie
& Win & Loss & Tie
& Win & Loss & Tie \\
\hline
\KHC L1   
& \heatmap{44} & \heatmap{44} & \heatmap{12}
& \heatmap{47} & \heatmap{38} & \heatmap{15}
& \heatmap{33} & \heatmap{53} & \heatmap{14} \\

\KHC LF   
& \heatmap{48} & \heatmap{45} & \heatmap{7}
& \heatmap{37} & \heatmap{34} & \heatmap{29}
& \heatmap{48} & \heatmap{44} & \heatmap{8} \\

\MHC L1  
& \heatmap{40} & \heatmap{42} & \heatmap{18}
& \heatmap{42} & \heatmap{38} & \heatmap{20}
& \heatmap{45} & \heatmap{48} & \heatmap{7} \\

\MHC LF  
& \heatmap{48} & \heatmap{41} & \heatmap{11}
& \heatmap{32} & \heatmap{33} & \heatmap{35}
& \heatmap{52} & \heatmap{40} & \heatmap{8} \\

\MRC L1 
& \heatmap{47} & \heatmap{46} & \heatmap{7}
& \heatmap{45} & \heatmap{29} & \heatmap{24}
& \heatmap{53} & \heatmap{35} & \heatmap{12} \\

\MRC LF 
& \heatmap{56} & \heatmap{31} & \heatmap{13}
& \heatmap{34} & \heatmap{30} & \heatmap{36}
& \heatmap{54} & \heatmap{38} & \heatmap{8} \\
\hline

\end{tabular}
}
\end{table}





\begin{table*}[!htb]
\centering
\caption{\threepointfive post-cutoff: Head-to-head win rates (\%), for comprehensiveness and diversity metrics. C0 and C1 are Leiden community levels from Edge et al.~\cite{edge24}. LF indicates leaf-level communities, and L1 indicates the level immediately above the leaf.}
\label{tab:gpt3_c0_c1}
\resizebox{\textwidth}{!}{%
\begin{tabular}{|r|ccc|ccc|ccc|ccc|ccc|ccc|}
\hline
\threepointfive
& \multicolumn{6}{c|}{\podcast} 
& \multicolumn{6}{c|}{\news} 
& \multicolumn{6}{c|}{\semi} \\
\cline{2-19}
results for
& \multicolumn{3}{c|}{C0} & \multicolumn{3}{c|}{C1}
& \multicolumn{3}{c|}{C0} & \multicolumn{3}{c|}{C1}
& \multicolumn{3}{c|}{C0} & \multicolumn{3}{c|}{C1} \\
\cline{2-19}
{\bf Comprehensiveness}
& Win & Loss & Tie & Win & Loss & Tie
& Win & Loss & Tie & Win & Loss & Tie
& Win & Loss & Tie & Win & Loss & Tie \\
\hline
\KHC L1   
& \heatmap{50} & \heatmap{42} & \heatmap{8} & \heatmap{68} & \heatmap{24} & \heatmap{8}
& \heatmap{56} & \heatmap{32} & \heatmap{12} & \heatmap{48} & \heatmap{42} & \heatmap{10}
& \heatmap{66} & \heatmap{30} & \heatmap{4} & \heatmap{72} & \heatmap{26} & \heatmap{2} \\

\KHC LF   
& \heatmap{58} & \heatmap{42} & \heatmap{0} & \heatmap{52} & \heatmap{44} & \heatmap{4}
& \heatmap{40} & \heatmap{50} & \heatmap{10} & \heatmap{42} & \heatmap{44} & \heatmap{14}
& \heatmap{68} & \heatmap{32} & \heatmap{0} & \heatmap{52} & \heatmap{46} & \heatmap{2} \\

\MHC L1  
& \heatmap{48} & \heatmap{46} & \heatmap{6} & \heatmap{66} & \heatmap{28} & \heatmap{6}
& \heatmap{50} & \heatmap{40} & \heatmap{10} & \heatmap{50} & \heatmap{36} & \heatmap{14}
& \heatmap{54} & \heatmap{46} & \heatmap{0} & \heatmap{62} & \heatmap{28} & \heatmap{10} \\

\MHC LF
& \heatmap{58} & \heatmap{40} & \heatmap{2} & \heatmap{68} & \heatmap{28} & \heatmap{4}
& \heatmap{44} & \heatmap{48} & \heatmap{8} & \heatmap{44} & \heatmap{49} & \heatmap{7}
& \heatmap{60} & \heatmap{38} & \heatmap{2} & \heatmap{72} & \heatmap{24} & \heatmap{4} \\

\MRC L1 
& \heatmap{58} & \heatmap{36} & \heatmap{6} & \heatmap{48} & \heatmap{48} & \heatmap{4}
& \heatmap{50} & \heatmap{38} & \heatmap{12} & \heatmap{46} & \heatmap{44} & \heatmap{10}
& \heatmap{54} & \heatmap{40} & \heatmap{6} & \heatmap{64} & \heatmap{32} & \heatmap{4} \\

\MRC LF 
& \heatmap{64} & \heatmap{36} & \heatmap{0} & \heatmap{57} & \heatmap{39} & \heatmap{4}
& \heatmap{56} & \heatmap{34} & \heatmap{10} & \heatmap{61} & \heatmap{35} & \heatmap{4}
& \heatmap{48} & \heatmap{42} & \heatmap{10} & \heatmap{55} & \heatmap{43} & \heatmap{2} \\
\hline
{\bf Diversity}
& Win & Loss & Tie & Win & Loss & Tie
& Win & Loss & Tie & Win & Loss & Tie
& Win & Loss & Tie & Win & Loss & Tie \\
\hline
\KHC L1   
& \heatmap{68} & \heatmap{28} & \heatmap{4} & \heatmap{60} & \heatmap{36} & \heatmap{4}
& \heatmap{60} & \heatmap{38} & \heatmap{2} & \heatmap{56} & \heatmap{38} & \heatmap{6}
& \heatmap{62} & \heatmap{38} & \heatmap{0} & \heatmap{66} & \heatmap{34} & \heatmap{0} \\

\KHC LF   
& \heatmap{52} & \heatmap{46} & \heatmap{2} & \heatmap{56} & \heatmap{40} & \heatmap{4}
& \heatmap{48} & \heatmap{48} & \heatmap{4} & \heatmap{50} & \heatmap{44} & \heatmap{6}
& \heatmap{62} & \heatmap{36} & \heatmap{2} & \heatmap{60} & \heatmap{34} & \heatmap{6} \\

\MHC L1  
& \heatmap{42} & \heatmap{54} & \heatmap{4} & \heatmap{70} & \heatmap{28} & \heatmap{2}
& \heatmap{56} & \heatmap{40} & \heatmap{4} & \heatmap{60} & \heatmap{34} & \heatmap{6}
& \heatmap{64} & \heatmap{36} & \heatmap{0} & \heatmap{70} & \heatmap{30} & \heatmap{0} \\

\MHC LF
& \heatmap{52} & \heatmap{48} & \heatmap{0} & \heatmap{72} & \heatmap{27} & \heatmap{1}
& \heatmap{60} & \heatmap{38} & \heatmap{2} & \heatmap{67} & \heatmap{32} & \heatmap{1}
& \heatmap{62} & \heatmap{38} & \heatmap{0} & \heatmap{52} & \heatmap{38} & \heatmap{10} \\

\MRC L1 
& \heatmap{66} & \heatmap{34} & \heatmap{0} & \heatmap{60} & \heatmap{40} & \heatmap{0}
& \heatmap{50} & \heatmap{48} & \heatmap{2} & \heatmap{42} & \heatmap{52} & \heatmap{6}
& \heatmap{52} & \heatmap{44} & \heatmap{4} & \heatmap{52} & \heatmap{46} & \heatmap{2} \\

\MRC LF 
& \heatmap{60} & \heatmap{40} & \heatmap{0} & \heatmap{65} & \heatmap{35} & \heatmap{0}
& \heatmap{64} & \heatmap{36} & \heatmap{0} & \heatmap{64} & \heatmap{36} & \heatmap{0}
& \heatmap{60} & \heatmap{40} & \heatmap{0} & \heatmap{56} & \heatmap{44} & \heatmap{0} \\
\hline
\end{tabular}
}
\end{table*}





\begin{table*}[!h]
\centering
\caption{\four post-cutoff : Head-to-head win rates (\%), for comprehensiveness and diversity metrics. C2 and C3 are Leiden community levels from Edge et al.~\cite{edge24}. LF indicates leaf-level communities, and L1 indicates the level immediately above the leaf.
}
\label{tab:gpt4}
\resizebox{\textwidth}{!}{%
\begin{tabular}{|r|ccc|ccc|ccc|ccc|ccc|ccc|}
\hline
\four
& \multicolumn{6}{c|}{\podcast} 
& \multicolumn{6}{c|}{\news} 
& \multicolumn{6}{c|}{\semi} \\
\cline{2-19}
results for
& \multicolumn{3}{c|}{C2} & \multicolumn{3}{c|}{C3}
& \multicolumn{3}{c|}{C2} & \multicolumn{3}{c|}{C3}
& \multicolumn{3}{c|}{C2} & \multicolumn{3}{c|}{C3} \\
\cline{2-19}
{\bf Comprehensiveness}
& Win & Loss & Tie & Win & Loss & Tie
& Win & Loss & Tie & Win & Loss & Tie
& Win & Loss & Tie & Win & Loss & Tie \\
\hline
\KHC L1   
& \heatmap{50} & \heatmap{50} & \heatmap{0} & \heatmap{50} & \heatmap{50} & \heatmap{0}
& \heatmap{38} & \heatmap{42} & \heatmap{20} & \heatmap{40} & \heatmap{42} & \heatmap{18}
& \heatmap{46} & \heatmap{46} & \heatmap{8} & \heatmap{39} & \heatmap{36} & \heatmap{25} \\

\KHC LF   
& \heatmap{50} & \heatmap{48} & \heatmap{2} & \heatmap{46} & \heatmap{48} & \heatmap{6}
& \heatmap{40} & \heatmap{40} & \heatmap{20} & \heatmap{43} & \heatmap{40} & \heatmap{17}
& \heatmap{52} & \heatmap{41} & \heatmap{7} & \heatmap{49} & \heatmap{32} & \heatmap{19} \\

\MHC L1  
& \heatmap{48} & \heatmap{46} & \heatmap{6} & \heatmap{48} & \heatmap{44} & \heatmap{8}
& \heatmap{42} & \heatmap{42} & \heatmap{16} & \heatmap{44} & \heatmap{44} & \heatmap{12}
& \heatmap{44} & \heatmap{40} & \heatmap{16} & \heatmap{51} & \heatmap{30} & \heatmap{19} \\

\MHC LF  
& \heatmap{50} & \heatmap{48} & \heatmap{2} & \heatmap{50} & \heatmap{46} & \heatmap{4}
& \heatmap{44} & \heatmap{40} & \heatmap{16} & \heatmap{46} & \heatmap{42} & \heatmap{12}

& \heatmap{52} & \heatmap{38} & \heatmap{10} & \heatmap{50} & \heatmap{24} & \heatmap{26} \\

\MRC L1 
& \heatmap{48} & \heatmap{44} & \heatmap{8} & \heatmap{44} & \heatmap{42} & \heatmap{14}
& \heatmap{38} & \heatmap{36} & \heatmap{26} & \heatmap{38} & \heatmap{32} & \heatmap{30}
& \heatmap{36} & \heatmap{42} & \heatmap{22} & \heatmap{46} & \heatmap{44} & \heatmap{10} \\

\MRC LF 
& \heatmap{48} & \heatmap{46} & \heatmap{6} & \heatmap{46} & \heatmap{48} & \heatmap{6}
& \heatmap{40} & \heatmap{35} & \heatmap{25} & \heatmap{40} & \heatmap{30} & \heatmap{30}
& \heatmap{48} & \heatmap{42} & \heatmap{10} & \heatmap{44} & \heatmap{44} & \heatmap{12} \\
\hline
{\bf Diversity}
& Win & Loss & Tie & Win & Loss & Tie
& Win & Loss & Tie & Win & Loss & Tie
& Win & Loss & Tie & Win & Loss & Tie \\
\hline
\KHC L1   
& \heatmap{50} & \heatmap{48} & \heatmap{2} & \heatmap{46} & \heatmap{48} & \heatmap{6}
& \heatmap{47} & \heatmap{40} & \heatmap{13} & \heatmap{52} & \heatmap{39} & \heatmap{9}
& \heatmap{38} & \heatmap{54} & \heatmap{8} & \heatmap{45} & \heatmap{48} & \heatmap{7} \\

\KHC LF   
& \heatmap{50} & \heatmap{50} & \heatmap{0} & \heatmap{48} & \heatmap{50} & \heatmap{2}
& \heatmap{50} & \heatmap{38} & \heatmap{12} & \heatmap{55} & \heatmap{37} & \heatmap{8}
& \heatmap{47} & \heatmap{49} & \heatmap{4} & \heatmap{52} & \heatmap{36} & \heatmap{12} \\

\MHC L1  
& \heatmap{50} & \heatmap{50} & \heatmap{0} & \heatmap{50} & \heatmap{50} & \heatmap{0}
& \heatmap{42} & \heatmap{48} & \heatmap{10} & \heatmap{44} & \heatmap{46} & \heatmap{10}
& \heatmap{50} & \heatmap{38} & \heatmap{12} & \heatmap{49} & \heatmap{36} & \heatmap{15} \\


\MHC LF  
& \heatmap{49} & \heatmap{46} & \heatmap{5} & \heatmap{46} & \heatmap{48} & \heatmap{6}
& \heatmap{44} & \heatmap{46} & \heatmap{10} & \heatmap{46} & \heatmap{44} & \heatmap{10}
& \heatmap{64} & \heatmap{32} & \heatmap{4} & \heatmap{58} & \heatmap{34} & \heatmap{8} \\

\MRC L1 
& \heatmap{48} & \heatmap{48} & \heatmap{4} & \heatmap{48} & \heatmap{44} & \heatmap{8}
& \heatmap{43} & \heatmap{34} & \heatmap{23} & \heatmap{46} & \heatmap{38} & \heatmap{16}
& \heatmap{46} & \heatmap{46} & \heatmap{8} & \heatmap{45} & \heatmap{45} & \heatmap{10} \\

\MRC LF 
& \heatmap{49} & \heatmap{50} & \heatmap{1} & \heatmap{48} & \heatmap{52} & \heatmap{0}
& \heatmap{45} & \heatmap{32} & \heatmap{23} & \heatmap{48} & \heatmap{36} & \heatmap{16}
& \heatmap{54} & \heatmap{44} & \heatmap{2} & \heatmap{54} & \heatmap{45} & \heatmap{1} \\
\hline
\end{tabular}
}
\end{table*}





\begin{table*}[!h]
\centering
\caption{\threepointfive post-cutoff: Head-to-head win rates (\%), for empowerment and directness metrics. C2/C3 are Leiden community levels, LF is leaf-level communities, and L1 is the level immediately above the leaf}
\label{tab:gpt35_empowerment_directness}
\resizebox{\textwidth}{!}{%
\begin{tabular}{|r|ccc|ccc|ccc|ccc|ccc|ccc|}
\hline
\threepointfive
& \multicolumn{6}{c|}{\podcast} 
& \multicolumn{6}{c|}{\news} 
& \multicolumn{6}{c|}{\semi} \\
\cline{2-19}
results for
& \multicolumn{3}{c|}{C2 } & \multicolumn{3}{c|}{C3 }
& \multicolumn{3}{c|}{C2 } & \multicolumn{3}{c|}{C3 }
& \multicolumn{3}{c|}{C2 } & \multicolumn{3}{c|}{C3 } \\
\cline{2-19}
{\bf Comprehensiveness}
& Win & Loss & Tie & Win & Loss & Tie
& Win & Loss & Tie & Win & Loss & Tie
& Win & Loss & Tie & Win & Loss & Tie \\
\hline
\KHC L1   
& \heatmap{48} & \heatmap{50} & \heatmap{2} & \heatmap{52} & \heatmap{45} & \heatmap{3}
& \heatmap{46} & \heatmap{52} & \heatmap{2} & \heatmap{63} & \heatmap{37} & \heatmap{0}
& \heatmap{52} & \heatmap{48} & \heatmap{0} & \heatmap{47} & \heatmap{52} & \heatmap{1} \\

\KHC LF   
& \heatmap{52} & \heatmap{43} & \heatmap{5} & \heatmap{54} & \heatmap{43} & \heatmap{3}
& \heatmap{46} & \heatmap{50} & \heatmap{4} & \heatmap{57} & \heatmap{40} & \heatmap{3}
& \heatmap{50} & \heatmap{50} & \heatmap{0} & \heatmap{50} & \heatmap{50} & \heatmap{0} \\

\MHC L1  
& \heatmap{54} & \heatmap{44} & \heatmap{2} & \heatmap{59} & \heatmap{39} & \heatmap{2}
& \heatmap{63} & \heatmap{35} & \heatmap{2} & \heatmap{55} & \heatmap{45} & \heatmap{0}
& \heatmap{40} & \heatmap{59} & \heatmap{1} & \heatmap{53} & \heatmap{43} & \heatmap{4} \\

\MHC LF  
& \heatmap{50} & \heatmap{46} & \heatmap{4} & \heatmap{54} & \heatmap{46} & \heatmap{0}
& \heatmap{46} & \heatmap{51} & \heatmap{3} & \heatmap{46} & \heatmap{54} & \heatmap{0}
& \heatmap{61} & \heatmap{38} & \heatmap{1} & \heatmap{52} & \heatmap{44} & \heatmap{4} \\

\MRC L1 
& \heatmap{39} & \heatmap{54} & \heatmap{7} & \heatmap{40} & \heatmap{55} & \heatmap{5}
& \heatmap{38} & \heatmap{60} & \heatmap{2} & \heatmap{38} & \heatmap{60} & \heatmap{2}
& \heatmap{51} & \heatmap{48} & \heatmap{2} & \heatmap{57} & \heatmap{37} & \heatmap{6} \\

\MRC LF 
& \heatmap{36} & \heatmap{56} & \heatmap{8} & \heatmap{47} & \heatmap{50} & \heatmap{3}
& \heatmap{50} & \heatmap{40} & \heatmap{10} & \heatmap{65} & \heatmap{30} & \heatmap{5}
& \heatmap{53} & \heatmap{46} & \heatmap{1} & \heatmap{59} & \heatmap{36} & \heatmap{5} \\
\hline
{\bf Directness}
& Win & Loss & Tie & Win & Loss & Tie
& Win & Loss & Tie & Win & Loss & Tie
& Win & Loss & Tie & Win & Loss & Tie \\
\hline
\KHC L1   
& \heatmap{62} & \heatmap{36} & \heatmap{2} & \heatmap{40} & \heatmap{50} & \heatmap{10}
& \heatmap{37} & \heatmap{63} & \heatmap{0} & \heatmap{45} & \heatmap{45} & \heatmap{10}
& \heatmap{35} & \heatmap{65} & \heatmap{0} & \heatmap{56} & \heatmap{43} & \heatmap{1} \\

\KHC LF   
& \heatmap{42} & \heatmap{55} & \heatmap{3} & \heatmap{35} & \heatmap{55} & \heatmap{10}
& \heatmap{40} & \heatmap{57} & \heatmap{3} & \heatmap{48} & \heatmap{42} & \heatmap{10}
& \heatmap{35} & \heatmap{65} & \heatmap{0} & \heatmap{58} & \heatmap{41} & \heatmap{1} \\

\MHC L1  
& \heatmap{54} & \heatmap{44} & \heatmap{2} & \heatmap{64} & \heatmap{28} & \heatmap{8}
& \heatmap{56} & \heatmap{44} & \heatmap{0} & \heatmap{56} & \heatmap{31} & \heatmap{13}
& \heatmap{63} & \heatmap{36} & \heatmap{1} & \heatmap{53} & \heatmap{43} & \heatmap{4} \\

\MHC LF  
& \heatmap{52} & \heatmap{44} & \heatmap{4} & \heatmap{62} & \heatmap{31} & \heatmap{7}
& \heatmap{46} & \heatmap{54} & \heatmap{0} & \heatmap{58} & \heatmap{31} & \heatmap{11}
& \heatmap{61} & \heatmap{38} & \heatmap{1} & \heatmap{53} & \heatmap{44} & \heatmap{3} \\

\MRC L1 
& \heatmap{43} & \heatmap{50} & \heatmap{7} & \heatmap{50} & \heatmap{44} & \heatmap{6}
& \heatmap{48} & \heatmap{50} & \heatmap{2} & \heatmap{59} & \heatmap{28} & \heatmap{13}
& \heatmap{50} & \heatmap{50} & \heatmap{0} & \heatmap{47} & \heatmap{53} & \heatmap{0} \\

\MRC LF 
& \heatmap{58} & \heatmap{36} & \heatmap{6} & \heatmap{47} & \heatmap{50} & \heatmap{3}
& \heatmap{48} & \heatmap{52} & \heatmap{0} & \heatmap{61} & \heatmap{28} & \heatmap{11}
& \heatmap{50} & \heatmap{50} & \heatmap{0} & \heatmap{45} & \heatmap{55} & \heatmap{0} \\
\hline
\end{tabular}
}
\vspace{-10ex}
\end{table*}



